\scanpage{66}

propres de la partie homogène de \deleted{\unsure{$u$ ou $u^{-1}$}} \inserted{l'un} quelconque ($u$ ou $u^{-1}$) des deux gén privilégiés de $G^+$, les axes de coordonnées sont les droites propres \deleted{pour} \inserted{de} $u$ (associé à l'épinglage) pour les val. propres $\zeta$ et $\zeta^{-1}$ resp\textsuperscript{t}.

\begin{center}\rule{1.5cm}{0.4pt}\end{center}

\uline{Formes} \deleted{\uline{quadratiques}} \uline{invariantes} (sous $\operatorname{Aut}(\Pi_{\zeta,k})$)

Si on regarde les fonctions \unsure{polygôniales|polygônales|polynômiales} \deleted{\illegible{1}} \inserted{sur} $E$ invariantes par $G$, comme l'origine est fixe, elles se décomposent en fonctions polygôniales homogènes, \deleted{En degré $m$} ou \uline{formes}

L'invariance par $\sigma$ signifie la symétrie p.r. aux \deleted{$x,y$} coordonnées $x,y$, donc

\[\left\{\begin{aligned}
&F(x,y)=G(\sigma_1(x,y),\sigma_2(x,y))\\
&\sigma_1(x,y)=x+y,\quad \sigma_2(x\unsure{.|,}y)=xy
\end{aligned}\right.\]

D'ailleurs, si \unsure{;|,|{}}

\[\deleted{\unsure{\operatorname{ord}}}F(x,y)=\sum_{\alpha+\beta=m}c_{\alpha\beta}\,x^\alpha y^\beta\]

on a

\[F(u(x,y))=F(\zeta x,\zeta^{-1}y)=\sum c_{\alpha\beta}\,\zeta^{\deleted{\unsure{\beta}}\alpha-\beta}x^\alpha y^{\unsure{\beta|\beta\zeta}}\]

donc on a

\[F\equiv F\circ u\Longleftrightarrow c_{\alpha\beta}=0\quad\text{si}\quad\beta-\alpha\not\equiv0\ (n)\]

% archive sheet number: 65
